\documentclass[11pt]{article}
\usepackage[letterpaper,hmargin=1in,vmargin=.87in]{geometry}
\usepackage{amsmath,amssymb,amsthm}
\usepackage{xcolor}
\definecolor{refblue}{HTML}{244C72}
\usepackage[colorlinks=true,linkcolor=black,citecolor=black,urlcolor=refblue,filecolor=refblue,
  pdftitle={A stretched-exponential lower bound for syntactically multilinear circuits}]{hyperref}
\newtheorem{theorem}{Theorem}
\newtheorem{corollary}[theorem]{Corollary}
\newtheorem{lemma}[theorem]{Lemma}
\newtheorem{proposition}[theorem]{Proposition}
\newcommand{\F}{\mathbb F}
\newcommand{\E}{\mathbb E}
\newcommand{\HC}{\mathrm{HC}}
\newcommand{\price}{\mathcal B}
\newcommand{\one}{\mathbf 1}
\newcommand{\ip}[2]{\langle #1,#2\rangle}
\newcommand{\norm}[1]{\lVert #1\rVert}
\setlength{\emergencystretch}{1.5em}
\raggedbottom
\title{A stretched-exponential lower bound for\\syntactically multilinear circuits}
\author{}
\date{}
\begin{document}
\maketitle
\vspace{-2.5em}
\begin{abstract}
We exhibit a polynomial-time Boolean function on \(N\) bits whose multilinear
representative over any field of characteristic two requires \(2^{N^\delta}\)
wires in any syntactically multilinear arithmetic circuit, for a fixed
\(\delta>0\). The bound holds for circuits of arbitrary depth with unbounded
fan-out. Deterministic decomposable circuits with internal negations, and
read-once parity branching programs with path-dependent query orders, can both
be converted with linear overhead into syntactically multilinear circuits for
the representative, so the same bound applies to them.
\end{abstract}

\section{One polynomial and two Boolean models}
Let \(d\ge4\) be even, \(k=d/2\), and \(N=d^2+d+1\). The inputs are a
bit \(b\), a vector \(x\in\{0,1\}^d\), and a matrix \(Y\in\{0,1\}^{d\times d}\).
Write \([d]=\{1,\ldots,d\}\) and \(I(x)=\{i:x_i=1\}\). Define
\[
 f_d(b,x,Y)=
 \begin{cases}
 \one\{Y\text{ is the permutation matrix of a single }d\text{-cycle}\},
       & b=0,\\
 \det_{\F_2}Y_{[k],I(x)}, & b=1,\ |I(x)|=k,\\
 0, & b=1,\ |I(x)|\ne k .
 \end{cases}
\]
The function \(f_d\) is computable in polynomial time: for \(b=0\), check that
\(Y\) is a permutation matrix and follow the orbit of 1; for \(b=1\), run
Gaussian elimination on the selected columns.
Let \(P_d\) be its unique multilinear representative over \(\F_2\), viewed
over any field \(F\) of characteristic two by scalar extension.
Uniqueness follows by interpolation in one Boolean variable at a time.

The \emph{scope} of an arithmetic gate is the set of variables labeling
inputs that reach it. A circuit is \emph{syntactically multilinear} if the
two input scopes of every multiplication are disjoint. Gates are binary
additions, subtractions, and multiplications; arbitrary field constants are
allowed; and size is the number of wires. A binary circuit with \(g\) arithmetic gates has \(2g\) input wires
to those gates, counted with multiplicity.

\begin{theorem}[Arithmetic lower bound]\label{thm:main}
There is a \(\delta>0\) such that for all sufficiently large even \(d\) and
every field \(F\) of characteristic two, any syntactically multilinear
arithmetic circuit over \(F\) computing \(P_d\) has at least \(2^{N^\delta}\)
wires.
\end{theorem}

\begin{corollary}[Deterministic decomposable circuits]\label{cor:dd}
Every deterministic decomposable circuit (d-D) computing \(f_d\), even with
arbitrary fan-out and negations at any gate, has at least
\(2^{N^{\Omega(1)}}\) wires.
\end{corollary}

\begin{corollary}[Read-once parity branching programs]\label{cor:bp}
Every syntactic read-once parity branching program computing \(f_d\)
has at least \(2^{N^{\Omega(1)}}\) edges, even if each path queries the
variables in its own order.
\end{corollary}

\paragraph{Relation to the general circuit bound.}
It is already known \cite[Theorem 1.1]{hc} that every arithmetic circuit having
the directed Hamiltonian-cycle polynomial \(\HC_d\) as its degree-\(d\)
homogeneous component needs
\(\exp(c(\log d)^2/(\log\log d)^2)\) gates, in every characteristic. By the
cycle restriction of Section~\ref{sec:cuts} and the two simulations of
Section~\ref{sec:sim}, that
quasi-polynomial bound applies to all three models above. The contribution
here is the improvement to \(2^{N^\delta}\). When characteristic two admits
ports, a property defined in Section~\ref{sec:proof}, this is
\cite[Theorem 12.2]{hc}. The new argument handles the case without ports,
where syntactic multilinearity lets us replace the low-rank cut summands of
\cite[Lemma 3.3]{hc} with exact products.

\section{Prior results}
\paragraph{Syntactically multilinear circuits.}
Raz \cite{raz} proved that multilinear formulas for the permanent and the
determinant have size \(n^{\Omega(\log n)}\), by bounding the rank of a
coefficient matrix across a partition of the variables. For syntactically
multilinear circuits, Raz, Shpilka, and Yehudayoff \cite{rsy} proved an
\(\Omega(n^{4/3}/\log^2 n)\) lower bound, and Alon, Kumar, and Volk \cite{akv}
improved it to \(\Omega(n^2/\log^2 n)\). Before \cite{hc} this was the best
bound known for the model. Lemma~\ref{lem:cut} uses the balanced-gate
argument of Valiant, Skyum, Berkowitz, and Rackoff \cite{vsbr}, which Raz and
Yehudayoff \cite{ry} showed preserves syntactic multilinearity.

\paragraph{Read-once parity branching programs.}
Jukna \cite{juknacodes} proved exponential lower bounds for oblivious
read-once parity branching programs. Brosenne, Homeister, and Waack \cite{bhw} and
Bollig, Waack, and Woelfel \cite{bww} did so for graph-driven programs, where
a single graph constrains the query order on every path. Without such a
constraint the best bound was the \(\Omega(n^2/\log^2 n)\) of Volk
\cite{volk}, and Jukna asked for an exponential bound
\cite[Research Problem 16.14]{juknabook}.

\paragraph{Deterministic decomposable circuits.}
Deterministic decomposable negation normal forms were introduced by Darwiche
\cite{darwiche} and studied in the knowledge compilation map of Darwiche and
Marquis \cite{dm}. Bova, Capelli, Mengel, and Slivovsky \cite{bcms} proved
exponential lower bounds for them through communication complexity. Those
bounds apply only to negation normal form, where negations occur only at the
inputs. Before \cite{hc}, we know of no superpolynomial lower bound for d-Ds
in which negations may occur at any gate.

\section{The two simulations}\label{sec:sim}
\begin{proof}[Proof of Corollary~\ref{cor:dd}]
In a d-D, the inputs to an AND have pairwise disjoint scopes. The inputs to
an OR are mutually exclusive on every Boolean assignment \cite[Definition 3]{dd}.
Binarizing a gate preserves both properties and increases the total wire
count by at most a constant factor. Over \(\F_2\), replace the three Boolean operations by
\[
 u\mathbin{\mathrm{AND}}v\longmapsto uv,\qquad
 u\mathbin{\mathrm{OR}}v\longmapsto u+v,\qquad
 \mathrm{NOT}\,u\longmapsto 1+u .
\]
At an OR, mutual exclusion makes the sum equal to the Boolean output. At an
AND, disjointness makes the multiplication syntactically multilinear.
An internal NOT introduces only a scalar and an addition. The translation
keeps every shared gate shared. By induction, every gate computes a
multilinear polynomial that agrees with the gate's Boolean value on the cube,
so the output equals \(P_d\). A d-D with \(s\) wires thus yields a
syntactically multilinear circuit with \(O(s+1)\) wires.
Theorem~\ref{thm:main} gives the bound.
\end{proof}

\begin{proof}[Proof of Corollary~\ref{cor:bp}]
A read-once parity branching program is an acyclic directed multigraph with
a source \(r\), a sink \(t\), and edges labeled by literals or Boolean
constants \cite[Definition 1.1]{volk}. Each variable is tested at most once
on each directed source--sink path. On a Boolean input the program outputs
the parity of the number of source--sink paths whose labels are all true.
The simulation below follows \cite[Claim 2.6]{volk}. Discard vertices and edges outside source--sink paths. Arithmetize a
positive literal \(z\) as \(z\), a negative literal as \(1+z\), and a
constant as itself. Write \(\ell_e\) for the resulting edge polynomial.
In topological order compute the prefix path sums
\[
       H_r=1,\qquad H_v=\sum_{e=(u,v)} H_u\,\ell_e\quad(v\ne r).
\]
The scope of \(H_u\) lies in the set of variables tested on source--\(u\)
paths. If \(e\) tests \(z\) and some source--\(u\) path also tests \(z\),
then that path followed by \(e\) and any \(v\)--\(t\) path tests \(z\)
twice, a contradiction. Thus the multiplication by \(\ell_e\) has disjoint input
scopes, even when prefixes merge after different query orders. The circuit
computes each \(H_u\) once. Binarizing its sums uses \(O(s+1)\) wires in
total for \(s\) graph edges. The value of \(H_t\) on the cube is the path
count modulo two, which is \(f_d\). Since \(H_t\) is multilinear, it equals
\(P_d\), and Theorem~\ref{thm:main} applies.
\end{proof}

\paragraph{Size conventions and padding.}
In both corollaries, taking any fixed positive exponent smaller than the one
in Theorem~\ref{thm:main} absorbs the constant-factor overhead of the
translation. For a d-D with \(g\) gates, merge identical input nodes and then
remove duplicate inputs to each gate by Boolean idempotence; at most
\((N+g+2)^2\) wires remain. For a read-once parity branching program with \(v\) vertices,
identically labeled parallel edges cancel in pairs, and at most
\((2N+2)v^2\) edges remain. Hence the same form of lower bound holds for the
number of gates of a d-D and the number of vertices of a read-once parity branching program. At
an arbitrary length \(n\), take the largest even \(d\) with \(N\le n\) and
ignore the remaining bits. Since \(N=\Theta(n)\) and pinning inputs to zero
preserves each of the three models, all three bounds hold at every
sufficiently large length.

\paragraph{An upper bound.}
The function \(f_d\) has a read-once parity branching program of size \(2^{O(d)}\).
The first edges test the selector \(b\). On the cycle branch, use states
\((S,v)\) with \(1\in S\subseteq[d]\) and \(v\in S\), recording a walk from
vertex 1 that has visited exactly \(S\) and is currently at \(v\). For each
\(u\notin S\), a chain of \(d\) edges from \((S,v)\) tests \(Y_{v,u}=1\) and
\(Y_{v,u'}=0\) for \(u'\ne u\), and leads to \((S\cup\{u\},u)\). From
\(([d],v)\) a final chain tests that row \(v\) is the indicator of column 1
and leads to the sink. A row is tested only when the walk leaves its
vertex, so every path is read-once. An accepting path tests every row;
hence one exists exactly when \(Y\) is the permutation matrix of a
\(d\)-cycle, and it is then unique. On the minor branch, use states \((i,R)\) with
\(0\le i\le d\) and \(R\subseteq[k]\), recording the rows matched to selected
columns among the first \(i\). From \((i-1,R)\) one edge tests \(x_i=0\) and
leads to \((i,R)\). Another tests \(x_i=1\), followed for each \(r\notin R\)
by an edge testing \(Y_{r,i}=1\) that leads to \((i,R\cup\{r\})\). The
accepting state is \((d,[k])\). Step \(i\) tests only \(x_i\) and entries of
column \(i\), so every path is read-once. The accepting paths correspond
to the bijections \(\sigma:I(x)\to[k]\) with all \(Y_{\sigma(i),i}=1\).
Their parity is the permanent of \(Y_{[k],I(x)}\), which equals the
determinant in characteristic two. Both branches have \(2^{d}d^{O(1)}\) edges. The
prefix simulation gives the same upper bound for syntactically multilinear
circuits, namely \(2^{O(\sqrt N)}\) in terms of the input length.

\section{Balanced cuts and the two restrictions}\label{sec:cuts}
The proof of Theorem~\ref{thm:main} uses two restrictions of the same
polynomial. With \(b=0\) we extract a homogeneous component. With \(b=1\)
we obtain a family of truth tables in \(x\), indexed by a matrix \(C\).
A \emph{balanced cut} is a set \(U\subseteq[d]\) with
\(d/3\le|U|\le2d/3\), together with its complement.
Lemma~\ref{lem:cut} shows that the circuit determines a single list of
balanced cuts that serves every table in the family.

\begin{lemma}[Balanced products under specialization]\label{lem:cut}
Let \(x_1,\ldots,x_d\), with \(d\ge4\), be distinguished variables of a
binary syntactically multilinear circuit with \(w\) wires. There is a list of balanced cuts \(U_1,\ldots,U_L\), with \(L\le w+1\),
such that every scalar specialization of the other variables gives an output
of the form
\[
                  \sum_{j=1}^{L} A_j(x_{U_j})B_j(x_{[d]\setminus U_j}).
\]
The factors may depend on the scalars, and the same list works for
specializations to any extension field.
\end{lemma}
\begin{proof}
Let \(S_g\) be the set of distinguished variables in the scope of \(g\),
before specialization, and let \(P_g\) be the polynomial computed at \(g\)
after specialization. Call \(g\) large if \(|S_g|>2d/3\) and small otherwise.
If the output is small, pad its scope to a balanced cut and take the other
factor to be 1.

For a large output, we maintain an expression for the output as a sum of
recorded cut products plus terms \(Q_gP_g\) over large gates \(g\), where
each multiplier \(Q_g\) uses only variables outside \(S_g\). Initially this
expression is \(1\cdot P_{\rm out}\). Process the large gates from the output
toward the inputs; before expanding \(g\), add the multipliers of all
occurrences of \(P_g\) into a single \(Q_g\).

At a sum or difference, pass \(Q_g\) to each child with its sign. For a
small child \(h\), record \((\pm Q_g)P_h\) on a cut with
\(S_h\subseteq U\subseteq S_g\) and
\(|U|=\max(|S_h|,\lceil d/3\rceil)\). Padding inside \(S_g\) leaves the
multiplier outside \(U\), even when \(h\) is a constant.

At a product, propagate to the child \(h\) with the larger set \(S_h\),
writing \(h'\) for its sibling. By disjointness, the new multiplier
\(Q_gP_{h'}\) uses no variable of \(S_h\), and \(|S_h|\ge |S_g|/2>d/3\).
If \(h\) is small, record \(P_h(Q_gP_{h'})\) across \(S_h\) and its complement.

Each expansion is an identity. Every input is small, and each wire
\((h,g)\) records at most one term, because all occurrences of \(P_g\) were
combined before \(g\) was expanded. Fixing ties and padding choices in advance and retaining
zero terms makes the list independent of the specialization.
\end{proof}

\paragraph{The cycle restriction.}
Set \(b=0\) and \(x=0\). Interpolation gives the formal identity
\[
 P_d(0,0,Y)=\sum_{\pi\text{ a }d\text{-cycle}}
      \prod_{i=1}^d Y_{i,\pi(i)}
      \prod_{(i,j):\,j\ne\pi(i)}(1+Y_{ij}).
\]
Each summand is the indicator that \(Y\) equals one particular permutation
matrix. Its degree-\(d\)
part is the product of the \(d\) positive literals, so the degree-\(d\)
homogeneous component is the directed Hamiltonian-cycle polynomial \(\HC_d\).

\paragraph{The minor restriction.}
Extend scalars from \(F\) to an uncountable algebraically closed field
\(K\) of characteristic two. Fix \(b=1\) and a Boolean vector \(x=\one_I\).
As a polynomial in \(Y\), the output equals the selected determinant when
\(|I|=k\), and zero otherwise. Indeed, both sides are multilinear in \(Y\)
and agree on every Boolean matrix. We can therefore substitute \emph{any}
\(C\in K^{k\times d}\) into the first \(k\) rows and zero into the rest.
The resulting truth table on the remaining \(d\) bits is
\[
 T_C(I)=
 \begin{cases}\det C_{[k],I},& |I|=k,\\0,& |I|\ne k.\end{cases}
\]
We call \(T_C\) a \emph{minor table}. Evaluating Lemma~\ref{lem:cut} on the \(x\)-cube expresses every \(T_C\)
as a sum of \(L\le w+1\) tables. Each summand factors across its listed
balanced cut. Crucially, the cuts are fixed \emph{before} \(C\) is chosen.

\section{Proof of the arithmetic lower bound}\label{sec:proof}
Throughout this section and Section~\ref{sec:bias} we use the terminology of
\cite{hc}: sites, suppliers, flags, sectors, legs, and prices.
A characteristic \emph{admits ports} if there are bounded-degree regular
graphs of unbounded order satisfying a common nonlinear Poincar\'e
inequality in all the spectral metric spaces of that characteristic
\cite[Section 2]{hc}. The
inequality bounds the mean squared distance of a random vertex pair by a
constant times that of a random edge, for every placement of the vertices.

\paragraph{If characteristic two admits ports.}
By \cite[Theorem 12.2]{hc}, there is an \(\varepsilon>0\) such that every
arithmetic circuit having \(\HC_d\) as its degree-\(d\) homogeneous component needs more
than \(2^{d^\varepsilon}\) gates for all large \(d\), over every field of
this characteristic. The cited bound even applies to limits of coefficient
vectors, as in border complexity. The cycle restriction exhibits \(\HC_d\)
as the degree-\(d\) homogeneous component of \(P_d(0,0,Y)\), so it gives the claimed
wire bound for any \(\delta<\varepsilon/2\).

\paragraph{If characteristic two does not admit ports.}
We use the minor restriction instead. The following proposition shows that
no short list of balanced cuts serves every minor table. Together with
Lemma~\ref{lem:cut}, it bounds the number of wires.

\begin{proposition}[Few cuts cannot serve every minor table]\label{prop:target}
Suppose characteristic two does not admit ports. Let
\(U_1,\ldots,U_L\) be balanced cuts, and suppose that for every \(C\in K^{k\times d}\) the table \(T_C\) is a sum of \(L\)
tables whose \(j\)th summand factors across \(U_j\) and its complement. Then
\begin{equation}\label{eq:target}
       \log(2L)=\Omega\!\left(\frac{\sqrt d}{(\log(2d))^{5/4}}\right).
\end{equation}
\end{proposition}

\paragraph{Weighted prices.}
We use the homogeneous spectral prices over \(K\), after adjoining Hilbert
coordinates \cite[Sections 2 and 3]{hc}.
Assign weight \(0\) to \(x_i=0\) and weight \(\theta_i\) to \(x_i=1\); we
call \(\theta_i\) the \emph{bias} of site \(i\). The
weighted price \(\price_\theta(T)\) is defined through tensor powers: fix the
total weight on each leg, add the sum of these weights to the sector's
ordinary price, maximize over sectors, and take the per-copy limit
\cite[Section 11]{hc}. Write
\[
       A(r,t)=r\log(1+e^{t/r}),\qquad S(\theta)=\sum_i A(1,\theta_i).
\]
Here \(S(\theta)\) is the transport budget for supplying all bits from the tip.

\paragraph{One bias per site.}
The proof of Proposition~\ref{prop:target} compares a lower bound on the
weighted price of one minor table with an upper bound on the price of every
cut product. The lower bound adapts \cite[Theorem 8.1]{hc}, which places
several columns at each site with evenly spaced biases. Here each site
carries a single bit \(x_i\) and gets a single column. Sampling the bias
together with each insertion makes one column per site suffice. We prove the
following lemma in Section~\ref{sec:bias}.

\begin{lemma}[One bias per site]\label{lem:bias}
Fix \(d\) even and unit sites \(u_1,\ldots,u_d\) with barycenter at the tip
of a metric cone carrying the prices above.
Let \(M=100\log(2d)\) and sample the \(\theta_i\) independently and
uniformly from the \(d^2\) midpoints of equal subintervals of \([-M,M]\).
With probability at least \(3/4\), there is a \(k\times d\) Vandermonde
matrix \(C\) over \(K\), where \(k=d/2\), for which
\(\price_\theta(T_C)\ge S(\theta)-B M\log d\), with \(B\) absolute.
\end{lemma}

\begin{proof}[Proof of Proposition~\ref{prop:target}]
Put \(m=\log(2d)\) and \(q=\log(2L)\), using natural logarithms.
The binary lamp certificate
\cite[Theorems 6.1 and 7.1]{hc} gives \(d\) unit sites \(u_i\) with
barycenter at the tip of a metric cone, and two suppliers per listed
cut. Both sides of each cut have at least \(d/3\) sites, which meets the
certificate's hypothesis that each side have at least \(d/4\).
Write \(z_{ji}\) for the supplier assigned to site \(i\) by cut \(j\).
With \(a_{ji}=(1+\norm{z_{ji}}^2-\operatorname{dist}(u_i,z_{ji})^2)/2\),
the certificate gives
\begin{equation}\label{eq:lamp}
  \sum_i a_{ji}\ge c d/\sqrt m,\qquad
  |a_{ji}|\le\norm{z_{ji}}\le R,\qquad R^2=O(m+q).
\end{equation}

\paragraph{A lower bound on the price of one minor table.}
Draw the biases \(\theta_i\) as in Lemma~\ref{lem:bias}, with
\(M=100m\). By that lemma, there is an absolute constant \(B\) such that
\begin{equation}\label{eq:near}
 \Pr_\theta\!\left[\exists\, C:\ \price_\theta(T_C)\ge
                    S(\theta)-B M\log d\right]\ge 3/4 .
\end{equation}
\paragraph{An upper bound on the price of each cut product.}
Place the suppliers for cut \(j\) at a fraction \(\tau\) of their displacement
from the tip. Each factor has zero ordinary price at its own supplier.
Sending bit \(i\) the distance \(r_{ji}=\operatorname{dist}(u_i,\tau z_{ji})\)
costs at most \(A(r_{ji},\theta_i)\). Sharing the summand index costs
\(O(\sqrt m\,q)\). Combining these with \cite[Lemmas 2.2 and 11.1]{hc}
gives, for \emph{every} \(C\),
\begin{equation}\label{eq:upper}
 \price_\theta(T_C)\le
  \max_{j\le L}\sum_i A(r_{ji},\theta_i)+O(\sqrt m\,q).
\end{equation}

\paragraph{Choosing the biases and the displacement.}
Suppose \(q\le\eta\sqrt d/m^{5/4}\) for a sufficiently small absolute
\(\eta>0\). We seek biases for which some minor table attains the lower
bound in \eqref{eq:near} and every listed cut gives a transport saving. Write
\(h(t)=\partial_r A(1,t)\) and \(\psi(t)=(1+t^2)e^{-|t|/2}\).
On \(1/2\le r\le2\) both \(|\partial_r A(r,t)|\) and
\(|\partial_r^2 A(r,t)|\) are bounded by \(C_0\psi(t)\). Indeed, the first
derivative is the binary entropy in nats at the logistic probability
\((1+e^{-t/r})^{-1}\), and the second is
\((t/r)^2 e^{t/r}/(r(1+e^{t/r})^2)\). The entropy is bounded below
near zero. The grid has spacing \(2M/d^2=o(1)\), and \(\psi,\psi^2\)
have integrable, eventually decreasing tails. Thus
\[
 \E h(\theta_i)\ge c/m,\qquad
 \E h(\theta_i)^2+\E\psi(\theta_i)\le C_1/m.
\]
We may take \(R\ge1\) in \eqref{eq:lamp}. For fixed \(j\), by independence
and \eqref{eq:lamp}, the sum \(\sum_i a_{ji}h(\theta_i)\) has mean at least
\(c_1d/m^{3/2}\), variance at most \(O(dR^2/m)\), and centered summands
bounded by \(O(R)\). Set \(t_0=c_1d/(2m^{3/2})\). Falling below \(t_0\)
requires a downward deviation of at least \(t_0\). Bernstein's inequality
bounds its probability by
\[
 \exp\!\left[-\Omega\!\left(\frac{t_0^2}{dR^2/m+Rt_0}\right)\right]
 \le \exp\!\left[-\Omega\!\left(\frac{d}{m^2(m+q)}\right)\right].
\]

Since \(m^2q(m+q)=o(d)\), the union of these bad events over \(L\)
cuts has probability \(o(1)\). Markov's inequality also gives
\(\sum_i\psi(\theta_i)\le C_2d/m\) with probability at least \(7/8\)
for a suitable \(C_2\). Since \((1-o(1))+7/8+3/4>2\), these two events and
the event in \eqref{eq:near} hold simultaneously for some biases. Fix such
biases and a matrix \(C\) witnessing \eqref{eq:near}.

Set \(\tau=m^{3/4}/\sqrt d\), so \(\tau R=o(1)\).
The cone law gives
\[
 r_{ji}^2=1-2\tau a_{ji}+\tau^2\norm{z_{ji}}^2,\qquad
 r_{ji}=1-\tau a_{ji}+O(\tau^2R^2).
\]
Taylor expansion on \(1/2\le r_{ji}\le2\) and the two simultaneous bounds
on the biases yield, for every \(j\),
\begin{align*}
 \sum_i A(r_{ji},\theta_i)
 &\le S(\theta)-\tau\sum_i a_{ji}h(\theta_i)
              +O\!\left(\tau^2R^2\sum_i\psi(\theta_i)\right)\\
 &\le S(\theta)-c_2\frac{\sqrt d}{m^{3/4}}
              +O\!\left(\sqrt m\,(m+q)\right).
\end{align*}
In \eqref{eq:upper}, the errors depending on \(q\) are
\(O(\sqrt m\,q)\le O(\eta\sqrt d/m^{3/4})\). Choose \(\eta\) small
enough to absorb them into half the saving. The remaining
\(O(m^{3/2})\) error and the allowed \(O(M\log d)\) deficit in
\eqref{eq:near} are both \(o(\sqrt d/m^{3/4})\). The lower bound
\eqref{eq:near} and the upper bound \eqref{eq:upper} on the price of this
\(T_C\) now contradict each other. Hence \(q>\eta\sqrt d/m^{5/4}\), which
is \eqref{eq:target}.
\end{proof}

\begin{proof}[Conclusion of Theorem~\ref{thm:main}]
Lemma~\ref{lem:cut} gives \(L\le w+1\). If characteristic two does not
admit ports, the minor
restriction meets the hypothesis of Proposition~\ref{prop:target}, so
\eqref{eq:target} and \(N=d^2+d+1\) imply
\[
        \log w=\Omega\!\left(\frac{N^{1/4}}{(\log N)^{5/4}}\right).
\]
Hence any \(0<\delta<1/4\) works in that case. If it admits ports, the
cycle restriction gives any \(0<\delta<\varepsilon/2\). Whether characteristic two
admits ports is a fixed fact independent of \(d\), and each case applies to
every sufficiently large even \(d\). Hence a single positive exponent works
for the whole family \(P_d\).
\end{proof}

\section{Proof of the bias lemma}\label{sec:bias}
We prove Lemma~\ref{lem:bias}. The exterior construction associates
threshold ranks with the columns of a matrix \cite[Sections 8--10]{hc}.
The argument follows the proof of \cite[Theorem 8.1]{hc}, with one column
per site and with random biases in place of the evenly spaced biases used
there at each site.

\paragraph{Centering an insertion.}
Each site \(u_i\) gives a flag \(F_i\) with standard logistic marginal.
In its pair law with a finite flag \(Q\), the pairing is the expectation
of the product. Since the barycenter is at the tip, the centering
inequality that closes the subsection on site flags and the barycenter in
\cite[Section 10]{hc} gives \(\sum_i\ip{Q}{F_i}\le0\).
The insertion sum of two flags, defined as twice their midpoint,
preserves addition of means and is 1-Lipschitz in each argument. Its
squared norm is the sum of the two squared norms plus twice the pairing
\cite[Section 9, insertion along a column curve]{hc}.
Pairing with \(Q\) is \(\norm{Q}\)-Lipschitz \cite[Proposition 9.3]{hc}.

Fix \(|t|\le M/2\), set \(c_t=1/2-t/(2M)\), and write
\(X(i,\xi)=\one_{\{F_i>t-\xi\}}-c_t\). For a uniform site \(I\) and
an independent grid bias \(\xi\),
\[
                  \E_{I,\xi}\ip{Q}{X(I,\xi)}
                       \le O(d^{-2})\norm{Q}.
\]
Indeed, at flag value \(f\), a continuous uniform bias on \([-M,M]\)
makes the mean indicator equal to \(c_t+f/(2M)\) clipped to \([0,1]\).
Clipping occurs only for \(|f|>M/2\). The logistic tails and \(M=100\log(2d)\)
give \(L^2\) clipping error \(O(d^{-10})\). Passing to the midpoint grid
changes the mean by at most \(d^{-2}\) pointwise. Apply Cauchy--Schwarz
in each pair law and average the remaining \(\ip{Q}{F_i}/(2M)\) over \(i\).

\paragraph{The second moment.}
Insert the sites in a uniform permutation \(\pi\) with independent grid
biases \(\xi_j\) indexed by position, so that \(\theta_{\pi(j)}=\xi_j\).
Reindexing independent biases by an independent permutation preserves their
joint law, so the \(\theta_i\) are distributed as in Lemma~\ref{lem:bias}.
Let \(Z_j\) be the iterated insertion sum of the first \(j\) centered
indicators. To bound the next pairing, swap the site at position \(j\)
with an independent uniform site \(I\), keeping position biases fixed.
The swapped permutation is uniform. Its earlier partial sum is within
distance two of \(Z_{j-1}\), since at most one earlier indicator changes,
each has norm at most one, and subsequent insertions are 1-Lipschitz.
The original \(Z_{j-1}\) is independent of \(I,\xi_j\). It follows that
\[
 \E\ip{Z_{j-1}}{X(\pi(j),\xi_j)}
       \le 2+O(d^{-2})\,\E\norm{Z_{j-1}}.
\]
Expanding the squared norm of the insertion and using
\(\norm{Z_{j-1}}\le j-1\) gives \(\E\norm{Z_d}^2=O(d)\), uniformly in \(t\).

\paragraph{From threshold ranks to price.}
Use the same order and biases at all thresholds
\(t=\pm M2^{-\ell}\), \(1\le \ell\le\lceil\log_2 d\rceil\).
By \cite[Lemma 9.4]{hc}, there is one \(k\times d\) Vandermonde matrix
\(C\) with distinct parameters for this finite list. Its threshold rank
\(r_t\), the dimension supplied by the columns at threshold \(t\), obeys
\[
    r_t\ge \E_{\rm flag}\min(k,V_t)-1,\qquad
    V_t=Z_d+dc_t,\qquad
    \E_{\rm flag}V_t=\sum_i p_i(t),\quad
    p_i(t)=(1+e^{t-\theta_i})^{-1}.
\]
Here \(\E_{\rm flag}\) averages over the flags with the realized biases and
order held fixed. Let \(\lambda(C)\) be the exterior value of \(C\). The
recovery formula at the start of \cite[Section 10]{hc} expresses it as
\(\int_0^\infty r_s\,ds-\int_{-\infty}^0(k-r_s)\,ds\). Hence the deficit
\(S(\theta)-\lambda(C)\) is the integral of \(K_+(t)=\sum_i p_i(t)-r_t\)
over \(t>0\) plus the integral of \(K_-(t)=k-r_t\) over \(t<0\).
Both losses are nonnegative and decrease away from zero, as in the proof
of \cite[Theorem 8.1]{hc}. The first is the kernel dimension of the column
map, whose domain shrinks as \(t\) grows. The second is the missing image
dimension. At a test
threshold the relevant loss is at most
\[
 1+\E_{\rm flag}\big(\operatorname{sgn}(t)Z_d-d|t|/(2M)\big)_+.
\]
Using \((z-v)_+\le z^2/(4v)\) for \(v>0\) and the second-moment bound
gives expected loss \(1+O(M/|t|)\) over the biases and order. On each
side of zero the innermost interval has length at most \(M/d\) and
loss at most \(d\). Each remaining dyadic interval up to \(M/2\) costs
\(O(M)\) by monotonicity, as does the interval from \(M/2\) to \(2M\)
in absolute value. Beyond \(2M\) use \(K_+\le\sum_i p_i\). Below
\(-2M\) any \(k\) basis columns supply independent dimensions, so
\(K_-\le\sum_{i\text{ in basis}}(1-p_i)\). The logistic tails cost
\(O(de^{-M})\) in total. Thus \(\E[S(\theta)-\lambda(C)]=O(M\log d)\).

In characteristic two, reordering wedge factors introduces no signs.
In the subset bases the exterior tensor on \(a\) copies is therefore
\(T_C^{\otimes a}\), since each coefficient is the product of the
selected minors over copies. Fixing the occupation count on each leg restricts
a weight sector to the corresponding exterior sector. Maximizing, per copy,
the price of that sector plus its weights gives
\(\price_\theta(T_C)\ge\lambda(C)\) by \cite[Sections 8 and 11]{hc}.
The exterior deficit is nonnegative by transport from the tip.
Markov's inequality at four times its expected bound gives success
probability at least \(3/4\). Since the existence of a suitable \(C\)
depends only on the biases \(\theta\), this proves Lemma~\ref{lem:bias}.

\begin{thebibliography}{15}
\bibitem{hc} OpenAI.
\href{https://github.com/openai/math/tree/main/preprints/Circuit-lower-bounds-for-Hamiltonian-cycle-and-the-permanent-November-24-2026}{\emph{Circuit lower bounds for Hamiltonian cycle and the permanent}}.
November 24, 2026 edition.
\bibitem{dd} M. Monet and D. Olteanu.
\href{https://www.cs.ox.ac.uk/dan.olteanu/papers/mo-amw18.pdf}{\emph{Towards deterministic decomposable circuits for safe queries}}.
AMW, 2018.
\bibitem{volk} B. L. Volk.
\href{https://arxiv.org/abs/2607.05944}{\emph{A lower bound for read-once parity branching programs}}. 2026.
\bibitem{raz} R. Raz.
\emph{Multi-linear formulas for permanent and determinant are of
super-polynomial size}. J. ACM 56(2), 2009.
\bibitem{rsy} R. Raz, A. Shpilka, and A. Yehudayoff.
\emph{A lower bound for the size of syntactically multilinear arithmetic
circuits}. SIAM J. Comput. 38(4):1624--1647, 2008.
\bibitem{akv} N. Alon, M. Kumar, and B. L. Volk.
\href{https://arxiv.org/abs/1708.02037}{\emph{Unbalancing sets and an almost
quadratic lower bound for syntactically multilinear arithmetic circuits}}.
Combinatorica 40(2):149--178, 2020.
\bibitem{vsbr} L. G. Valiant, S. Skyum, S. Berkowitz, and C. Rackoff.
\emph{Fast parallel computation of polynomials using few processors}.
SIAM J. Comput. 12(4):641--644, 1983.
\bibitem{ry} R. Raz and A. Yehudayoff.
\emph{Balancing syntactically multilinear arithmetic circuits}.
Comput. Complexity 17:515--535, 2008.
\bibitem{juknacodes} S. Jukna.
\emph{Linear codes are hard for oblivious read-once parity branching
programs}. Inform. Process. Lett. 69:267--269, 1999.
\bibitem{bhw} H. Brosenne, M. Homeister, and S. Waack.
\emph{Characterizing the complexity of Boolean functions represented by
well-structured graph-driven parity-FBDDs}.
RAIRO Theor. Inform. Appl. 36(3):229--253, 2002.
\bibitem{bww} B. Bollig, S. Waack, and P. Woelfel.
\emph{Parity graph-driven read-once branching programs and an exponential
lower bound for integer multiplication}. IFIP TCS, 2002.
\bibitem{juknabook} S. Jukna.
\emph{Boolean Function Complexity: Advances and Frontiers}. Springer, 2012.
\bibitem{darwiche} A. Darwiche.
\emph{On the tractable counting of theory models and its application to
truth maintenance and belief revision}.
J. Appl. Non-Classical Logics 11(1--2):11--34, 2001.
\bibitem{dm} A. Darwiche and P. Marquis.
\emph{A knowledge compilation map}.
J. Artificial Intelligence Res. 17:229--264, 2002.
\bibitem{bcms} S. Bova, F. Capelli, S. Mengel, and F. Slivovsky.
\emph{Knowledge compilation meets communication complexity}.
IJCAI, 1008--1013, 2016.
\end{thebibliography}
\end{document}
